\documentclass[12pt]{article}
\usepackage[utf8]{inputenc}
\usepackage[T1]{fontenc}
\usepackage{hyperref}
\title{Tecpatl\'an Totonac Language Resources for K-12 Curriculum Development}
\author{Compiled by Claude Sonnet 5 (claude-sonnet-5) for Andr\'as Kornai}
\date{}
\begin{document}
\maketitle

\section{Introduction}
Tecpatl\'an Totonac (ISO 639-3 tcw; Glottocode tecp1235) is a Totonacan language spoken in Puebla, Mexico. Wikipedia cites about 540 speakers (2000 census, via Ethnologue), while Joshua Project gives 700. The reference description says it is not known to be taught in schools. Search found very little: no dictionary, grammar or learning material.

\section{Primary Educational Resources}
None found.

\section{Linguistic and Reference Materials}
\begin{itemize}
\item Wikipedia article on Tecpatl\'an Totonac: \url{https://en.wikipedia.org/wiki/Tecpatl%C3%A1n_Totonac}
\item Glottolog language page (references list, including subclassification comments): \url{https://glottolog.org/resource/languoid/iso/tcw}
\item Wikipedia overview of the Totonacan languages: \url{https://en.wikipedia.org/wiki/Totonacan_languages}
\end{itemize}

\section{Religious and Community Resources}
\begin{itemize}
\item Joshua Project language page. It reports Bible portions (1988--1995) and a New Testament (1999), with no complete Bible: \url{https://joshuaproject.net/index.php/languages/tcw}
\item Joshua Project people-group page: \url{https://joshuaproject.net/people_groups/19673}
\end{itemize}

\section{Digital and Audio Resources}
Global Recordings Network lists no recordings in this language: \url{https://globalrecordings.net/en/language/tcw}

\section{Curriculum Development Context}
There is almost no published material. A curriculum would depend on new documentation with the community. The Glottolog references and the New Testament translation are the only leads found here.
\end{document}